\documentclass[10pt,slidestop]{beamer}
\synctex=1

\usepackage{etex}
\usepackage{fourier-orns}
\usepackage{ccicons}
\usepackage{amssymb}
\usepackage{amstext}
\usepackage{amsbsy}
\usepackage{amsopn}
\usepackage{amscd}
\usepackage{amsxtra}
\usepackage{amsthm}
\usepackage{float}
\usepackage{color}
\usepackage{mathrsfs}
\usepackage{bm}
\usepackage[nice]{nicefrac}
\usepackage{setspace}
\usepackage{hyperref}
\usepackage{ragged2e}
\usepackage{listings}
\usepackage{fontawesome}
\usepackage{environ}
\usepackage{manfnt}

%\usepackage[utf8x]{inputenc}
\usepackage[T1]{fontenc}
\usepackage{array}
\usepackage[super]{nth}
\usepackage{algorithm}
\usepackage{algpseudocode}
\newcolumntype{x}[1]{>{\centering\arraybackslash\hspace{0pt}}m{#1}}
\setlength{\parindent}{0pt}
\pdfcompresslevel=9

\usepackage{verbatimbox}
\usepackage{pgfplots}
\usepackage{matlab-prettifier}
\usepackage{listings}
\usepackage{xcolor}

\definecolor{eclipseBlue}{RGB}{42,0.0,255}
\definecolor{eclipseGreen}{RGB}{63,127,95}
\definecolor{eclipsePurple}{RGB}{127,0,85}
\definecolor{eclipseEmph}{RGB}{0,0,0}
\definecolor{macroOrange}{RGB}{176,79,0}

% MATLAB console sessions and generated scripts: MATLAB's keywords, comments and
% strings are coloured as in the editor; the prompt and the output stay plain.
\lstdefinestyle{MatlabConsole}{
  language=Matlab,
  basicstyle=\ttfamily\small,
  keywordstyle=\color{eclipsePurple}\bfseries,
  commentstyle=\color{eclipseGreen},
  stringstyle=\color{eclipseBlue},
  backgroundcolor=\color{black!5},
  frame=single,
  breaklines=true,
  numbers=none,
  showstringspaces=false,
}

\lstdefinelanguage{Dynare}
{
  % list of keywords
  morekeywords={
    var_model,
    trend_component_model,
    var_expectation_model,
    var_expectation,
    pac_model,
    pac_expectation,
    model,
    diff,
    end,
    varexo,
    var,
    parameters,
    steady_state_model,
    steady,
    shocks,
    perfect_foresight_setup,
    perfect_foresight_solver,
    initval,
    endval,
    histval,
    stoch_simul,
    check,
    resid,
    estimation,
    varobs,
    verbatim,
    predetermined_variables,
  },
  sensitive=false, % keywords are not case-sensitive
  morecomment=[l]{//}, % l is for line comment
  morecomment=[l]{\%},
  morecomment=[s]{/*}{*/}, % s is for start and end delimiter
  morestring=[b]',         % defines that strings are enclosed in single quotes
  morestring=[b]",
  % The macro language: a directive line from its @# to the end of the line, and a
  % substitution @{...} inside the text, in a colour of their own.
  moredelim=[l][\color{macroOrange}\bfseries]{@\#},
  moredelim=[s][\color{macroOrange}]{@\{}{\}},
}

% Set Language
\lstset{
  language={Dynare},
  inputencoding = utf8,
  basicstyle=\footnotesize\ttfamily, % Global Code Style
  captionpos=b, % Position of the Caption (t for top, b for bottom)
  extendedchars=true, % Allows 256 instead of 128 ASCII characters
  tabsize=2, % number of spaces indented when discovering a tab
  columns=fixed, % make all characters equal width
  keepspaces=true, % does not ignore spaces to fit width, convert tabs to spaces
  showstringspaces=false, % lets spaces in strings appear as real spaces
  breaklines=true, % wrap lines if they don't fit
  frame=trbl, % draw a frame at the top, right, left and bottom of the listing
  frameround=tttt, % make the frame round at all four corners
  framesep=4pt, % quarter circle size of the round corners
  numbers=left, % show line numbers at the left
  numberstyle=\tiny\ttfamily, % style of the line numbers
  commentstyle=\color{eclipseGreen}, % style of comments
  keywordstyle=\color{eclipsePurple}\bfseries, %, % style of keywords
  stringstyle=\color{eclipseBlue}, % style of strings
  emph={model_name, periods, eqtags, targets, expression, auxiliary_model_name, horizon,discount},emphstyle=\color{eclipseEmph}\bfseries,
  literate={⟂}{{\ensuremath{\perp}}}{1},
}



\newcommand{\trace}{\mathrm{tr}}
\newcommand{\vect}{\mathrm{vec}}
\newcommand{\tracarg}[1]{\mathrm{tr}\left\{#1\right\}}
\newcommand{\vectarg}[1]{\mathrm{vec}\left(#1\right)}
\newcommand{\vecth}[1]{\mathrm{vech}\left(#1\right)}
\newcommand{\iid}[2]{\mathrm{iid}\left(#1,#2\right)}
\newcommand{\normal}[2]{\mathcal N\left(#1,#2\right)}
\newcommand{\dynare}{\href{http://www.dynare.org}{\color{blue}Dynare}}
\newcommand{\AllSample}{ \mathcal Y_T }
\newcommand{\sample}{\mathcal Y_T}
\newcommand{\samplet}[1]{\mathcal Y_{#1}}
\newcommand{\slidetitle}[1]{\fancyhead[L]{\textsc{#1}}}

\definecolor{gray}{gray}{0.4}

\theoremstyle{plain}

\makeatletter
\@ifclassloaded{beamer}{
\setbeamertemplate{footline}{
  {\hfill\vspace*{1pt}\href{https://creativecommons.org/publicdomain/zero/1.0/}{\cczero}\hspace{.1cm}\today\enspace
  }}
\setbeamertemplate{navigation symbols}{}
\setbeamertemplate{blocks}[rounded][shadow=true]
\setbeamertemplate{caption}[numbered]}
\makeatother

% From https://tex.stackexchange.com/questions/312118/a-note-with-the-dangerous-bend-symbol-how-to-vertically-center-text-and-symbol
\newcommand{\dnote}[1]{%
    \noindent % I guess this is intended...
    \begin{tabular}{@{}m{0.13\textwidth}@{}m{0.87\textwidth}@{}}%
        \huge\textdbend &#1%
    \end{tabular}%
    \par % ... and this too.
}

\NewEnviron{notes}{\justifying\footnotesize\begin{spacing}{1.0}\BODY\vfill\pagebreak\end{spacing}}

\begin{document}

\title{modBuilder}
\author[]{\texttt{stephane.adjemian@univ-lemans.fr}}
\date{September, 2026}

\begin{frame}
  \titlepage{}
\end{frame}

\begin{frame}[fragile]
  \frametitle{Introduction}

  \bigskip

  \begin{itemize}

  \item \verb+modBuilder+ is a MATLAB class aimed at simplifying the interactive and
    programmatic creation of Dynare \verb+.mod+ files.\newline

  \item Build incrementally a model directly from the MATLAB.\newline

  \item Facilitates consistency checks, error detection, and modular
    model development.\newline

  \item Also provides tools for computing the steady state: solving equations
    one at a time, or deriving the whole steady state structurally.\newline

  \item Export the model to a \verb+.mod+ file ready to run, or to LaTeX.\newline

  \item Code available in a public \href{https://git.dynare.org/stepan-a/modbuilder}{Git repository} (not part of Dynare).

  \end{itemize}

\end{frame}


\begin{frame}[fragile]
  \frametitle{Instantiate an empty model}

\begin{lstlisting}[style=MatlabConsole]
>> model = modBuilder()

model =

  modBuilder with properties:

       params: {0x4 cell}
       varexo: {0x4 cell}
          var: {0x4 cell}
         tags: [1x1 struct]
      symbols: {1x0 cell}
    equations: {0x2 cell}
            T: [1x1 struct]
         date: 04-Oct-2025 09:39:05
   \end{lstlisting}

\end{frame}


\begin{frame}[fragile]
  \frametitle{Instantiate a model from a \texttt{.mod} file -- I --}

  \lstinputlisting[firstline=1,lastline=19,frame=none,basicstyle=\footnotesize,language=Dynare]{../models/rbc12.mod}

\end{frame}

\begin{frame}[fragile]
  \frametitle{Instantiate a model from a \texttt{.mod} file -- II --}

  \lstinputlisting[firstline=21,lastline=41,frame=none,basicstyle=\footnotesize,language=Dynare]{../models/rbc12.mod}

\end{frame}


\begin{frame}[fragile]
  \frametitle{Instantiate a model from a \texttt{.mod} file -- III --}

  \begin{lstlisting}[style=MatlabConsole]

    >> model = modBuilder('rbc12.mod')

model =

  modBuilder with properties:

       params: {8x4 cell}
       varexo: {2x4 cell}
          var: {6x4 cell}
         tags: [1x1 struct]
      symbols: {1x0 cell}
    equations: {6x2 cell}
            T: [1x1 struct]
         date: 04-Oct-2025 09:25:42
   \end{lstlisting}

  \begin{itemize}
  \item No Dynare run, no JSON file: the \texttt{.mod} file is read directly.
  \end{itemize}

\end{frame}


\begin{frame}[fragile]
  \frametitle{Instantiate a model from a \texttt{.mod} file -- III bis --}

  Reading a \texttt{.mod} file writes a \texttt{modBuilder} script and runs it. Keep
  the script and carry on from there:

  \begin{lstlisting}[style=MatlabConsole,basicstyle=\ttfamily\footnotesize]
    >> model = modBuilder('rbc12.mod', Script='rbc12.m');
   \end{lstlisting}

  \begin{lstlisting}[style=MatlabConsole,basicstyle=\ttfamily\footnotesize]
% modBuilder script generated from rbc12.mod. Edit freely.

m = modBuilder();

% Equations
m.add('a', 'a = rho*a(-1)+tau*b(-1)+e');
m.add('y', 'y = exp(a)*k(-1)^alpha*h^(1-alpha)');
...
% Calibration
alpha = 0.360000;
...
% Parameters
m.parameter('alpha', alpha, 'texname', '\alpha', 'declared', true);
  \end{lstlisting}

\end{frame}


\begin{frame}[fragile]
  \frametitle{Instantiate a model from a \texttt{.mod} file -- macros --}

  The macro language is read in full; its control flow survives into the script.

  \begin{lstlisting}[language=Dynare,basicstyle=\scriptsize\ttfamily]
@#define Countries = ["US", "EA"]
model;
@#for c in Countries
[name = 'y_@{c}']
y_@{c} = alpha + e_@{c};
@#endfor
@#if OpenEconomy
[name = 'nx']
nx = y_US - y_EA;
@#endif
end;
  \end{lstlisting}

  becomes

  \begin{lstlisting}[style=MatlabConsole,basicstyle=\ttfamily\footnotesize]
Countries = macroarray('US', 'EA');
m.add('y_$1', 'y_$1 = alpha + e_$1', Countries);
if OpenEconomy
    m.add('nx', 'nx = y_US - y_EA');
end
  \end{lstlisting}

\end{frame}


\begin{frame}[fragile]
  \frametitle{Instantiate a model from a \texttt{.mod} file -- macros II --}

  \begin{itemize}
  \item Covered: \texttt{@\#define} for variables and functions, with Dynare's own
    scoping; \texttt{@\#if}, \texttt{@\#ifdef}, \texttt{@\#ifndef};
    \texttt{@\#for} over tuples and with a \texttt{when} guard; \texttt{@\#include};
    comprehensions, casts, set operations and ranges; \texttt{@\#echo} and
    \texttt{@\#error}; continuation lines. Dynare's own macro-processor self-test
    loads end to end.
  \item A loop becomes one implicit-loop call when a \texttt{\$N} template reproduces
    every iteration, and a MATLAB \texttt{for} loop otherwise. The template is guessed
    and then verified against the iterations, so a wrong guess is discarded rather than
    emitted. The loop runs over the local its set became, \texttt{Countries} above, so
    editing that local changes the model; a \texttt{when} guard becomes a
    \texttt{filter}.
  \item A \texttt{@\#if} becomes a MATLAB \texttt{if}/\texttt{elseif}/\texttt{else}
    carrying every branch. A branch that is not selected produces no text, so the file is
    read again with the conditional forced onto it; the model itself always comes from
    the unforced read. Changing the flag in the script then gives the model the
    \texttt{.mod} file would have given with that setting.
  \end{itemize}

\end{frame}


\begin{frame}[fragile]
  \frametitle{Instantiate a model from a \texttt{.mod} file -- macros, \texttt{-D} --}

  A flag with a default, that Dynare's \texttt{-DOpen=true} overrides:

  \begin{lstlisting}[language=Dynare,basicstyle=\scriptsize\ttfamily]
@#ifndef Open
@#define Open = false
@#endif
...
@#if Open
c = y - e;
@#else
c = y;
@#endif
  \end{lstlisting}

  \texttt{Defines=} is \texttt{-D}. The flag is a local of the script, at the value the
  read used, and the \texttt{if} keeps both branches:

  \begin{lstlisting}[style=MatlabConsole,basicstyle=\ttfamily\scriptsize]
>> model = modBuilder('open.mod', Defines=struct('Open', true));

Open = true;
...
if Open
    m.add('c', 'c = y - e');
else
    m.add('c', 'c = y');
end
  \end{lstlisting}

\end{frame}


\begin{frame}[fragile]
  \frametitle{Instantiate a model from a \texttt{.mod} file -- what is read --}

  \begin{itemize}
  \item \textbf{Imported:} the declarations, the calibration, the \texttt{model} block
    with its tags and its \texttt{\#} model-local variables, \texttt{steady\_state\_model},
    \texttt{initval}, and the options of \texttt{steady} and \texttt{check}.\newline
  \item \textbf{Skipped, with a warning:} everything computational
    (\texttt{stoch\_simul}, \texttt{estimation}, \texttt{shocks}, \dots) and every native
    MATLAB line, told apart by Dynare's own rule: a line whose first token is neither a
    keyword nor a declared symbol is MATLAB. A native assignment that a calibration or an
    initial value uses is kept as a local of the script.
    \texttt{Strict=true} turns the warnings into errors.\newline
  \item \textbf{Refused}, because skipping them would silently build a different model:
    \texttt{predetermined\_variables}, \texttt{varexo\_det}, trend variables, external
    functions, the optimal-policy statements, a second \texttt{model} block.
  \end{itemize}

  \medskip
  \small Of the 280 files of Dynare's test suite that use the macro language, 222 load
  as they are. Most of the rest hit a refused statement or an operator the parser does
  not know yet (\texttt{$\perp$}); a few expect a macro variable that Dynare's test
  runner sets with \texttt{-D}, and load once \texttt{Defines=} sets it.

\end{frame}


\begin{frame}[fragile]
  \frametitle{Instantiate a model from a \texttt{.mod} file -- naming equations --}

  A \texttt{modBuilder} equation is keyed to one endogenous variable. A \texttt{name}
  tag that is a declared variable settles it. Without tags the reader pairs equations
  and variables itself (next frame), and when no pairing covers every variable, the
  error names what to tag. Here neither equation mentions \texttt{c}:

  \begin{lstlisting}[language=Dynare,basicstyle=\scriptsize\ttfamily]
var y c;
varexo e;
parameters alpha;
alpha = 0.5;

model;
y = alpha*e;
y = 2*e;
end;
  \end{lstlisting}

  \begin{lstlisting}[style=MatlabConsole,basicstyle=\ttfamily\scriptsize]
>> modBuilder('twoeqs.mod')
twoeqs.mod: unable to associate every equation with a unique
endogenous variable. Add a "name" tag to these:
  line 8: y = 2*e
  unmatched endogenous variables: c
  \end{lstlisting}

\end{frame}


\begin{frame}[fragile]
  \frametitle{Instantiate a model from a \texttt{.mod} file -- pairing --}

  Every equation of a \texttt{modBuilder} model is the equation \emph{of} one endogenous
  variable, and every variable has one. Without \texttt{name} tags the reader has to
  decide which is which: $n$ equations, $n$ variables, each equation to be given a
  distinct variable it can stand for.

  \medskip
  An equation may stand for a variable it mentions. Two cases, decided on the static
  equation, leads and lags dropped and the result simplified:
  \begin{enumerate}\small
  \item the static equation pins the variable: it holds it outright, as
    \texttt{y = alpha*k(-1) + e} holds \texttt{y} and \texttt{k}, or as a multiple whose
    other factor cannot vanish, as \texttt{junk = 0.9*junk(+1)}, which reduces to
    $0.1\,\texttt{junk} = 0$ and fixes \texttt{junk} at zero;
  \item the static equation leaves the variable free: an Euler equation reduces to
    $c^{-\sigma}(1 - \beta R) = 0$, which pins $R$ and says nothing about $c$; and
    \texttt{Y/Y(-1) = g} reduces to \texttt{1 = g}, with \texttt{Y} gone. Both are
    still the equations of \texttt{c} and of \texttt{Y}, only weaker claims.
  \end{enumerate}

  \medskip
  When an equation could stand for several variables, two preferences settle it: the
  variable on its left-hand side, if it is one of them; otherwise the variable that the
  fewest other equations could stand for, since a variable mentioned by one equation
  only must take that one or end up with none.

\end{frame}


\begin{frame}[fragile]
  \frametitle{Instantiate a model from a \texttt{.mod} file -- pairing, the choice --}

  Each possible pair is given a price: 1 when the static equation pins the variable, 2
  when it leaves it free, less a discount for the left-hand side, plus a small surcharge
  when many equations want the same variable. Leaving an
  equation without a variable is priced above any pair. The cheapest complete assignment
  is then computed exactly; it is the problem of assigning $n$ workers to $n$ jobs at the
  least total cost, and MATLAB solves it (\texttt{matchpairs}). Good pairs are used
  wherever they can be, and give way only where the assignment as a whole needs it.

  \begin{center}\small
  \begin{tabular}{l|cccc}
                                     & \texttt{y} & \texttt{k} & \texttt{i} & \texttt{junk} \\ \hline
    \texttt{y = alpha*k(-1) + e}     & \textbf{1} & 1          &            &               \\
    \texttt{k = (1-delta)*k(-1) + i} &            & \textbf{1} & 1          &               \\
    \texttt{i = s*y}                 & 1          &            & \textbf{1} &               \\
    \texttt{junk = 0.9*junk(+1)}     &            &            &            & \textbf{1}    \\
  \end{tabular}
  \end{center}
  \small The case of each possible pair, in bold the assignment chosen: the left-hand
  side settles the first three rows, and the last equation takes \texttt{junk}, which it
  pins at zero.

  When no complete assignment exists, some variable appears in no equation left for it,
  the reader stops and lists the equations to tag rather than guess. A \texttt{name} tag
  always wins over the assignment.

\end{frame}


\begin{frame}[c,fragile]
  \frametitle{Instantiate a model}

  \begin{itemize}

  \item It is also possible to create a new model from a previously defined modBuilder object that has been saved to disk.\newline

  \item Once a \verb+modBuilder+ object is instantiated, the model can be expanded or transformed by utilizing the methods provided by the class.\newline

  \item In the end the model can be saved on disk for later usage (\verb+save+ method) or exported to a \verb+.mod+ file (\verb+write+ method).

  \end{itemize}

\end{frame}


\begin{frame}[c,fragile]
  \frametitle{Add an equation}

  \begin{lstlisting}[style=MatlabConsole]

    >> model.add('y', 'y = rho*y(-1)-gamma*z+e');
   \end{lstlisting}

   \bigskip

  \begin{itemize}

  \item Each equation must be associated to an endogenous variable (\verb+y+) that will be used to index the model object.\newline

  \item[$\Rightarrow$] Each new equation adds a new endogenous variable to the model.\newline

  \item[$\Rightarrow$] Removing the equation of \verb+y+ therefore says what becomes of
    \verb+y+ (next frame).\newline

  \item All the other symbols (\verb+rho+, \verb+z+, and \verb+e+), unless previously defined, are classified as unresolved.\newline

  \item These symbols need to be typed before completing the model. This can be done by adding an equation for an endogenous variable (\verb+z+) or explictly declaring a symbol as a parameter (\verb+rho+) or an exogenous variable (\verb+e+).

  \end{itemize}

\end{frame}


\begin{frame}[fragile]
  \frametitle{Remove an equation}

  An equation is the equation \emph{of} a variable, so removing it decides the fate of
  that variable, and of the symbols only that equation used:

  \begin{lstlisting}[style=MatlabConsole,basicstyle=\ttfamily\footnotesize]
m.add('y', 'y = a*k(-1) + e');
m.add('k', 'k = (1-delta)*k(-1) + i');
m.add('i', 'i = s*y');
m.add('h', 'h = theta*y');
  \end{lstlisting}

  \begin{center}\footnotesize
  \begin{tabular}{@{}lll@{}}
    \hline
    call & what happens & endogenous \,/\, exogenous \\
    \hline
    \verb+m.remove('i')+ & \verb+i+ is still used by the equation of \verb+k+: & \verb+y k h+ \,/\, \verb+e i+ \\
                         & it becomes exogenous, \verb+s+ goes & \\
    \verb+m.remove('h')+ & \verb+h+ is used nowhere else: & \verb+y k i+ \,/\, \verb+e+ \\
                         & it goes, and \verb+theta+ with it & \\
    \verb+m.rmflip('i', 'k')+ & \verb+i+ stays endogenous, keyed to & \verb+y i h+ \,/\, \verb+e k+ \\
                         & the former equation of \verb+k+ & \\
    \hline
  \end{tabular}
  \end{center}

  \medskip
  \small Which variable an equation belongs to is therefore not a label: it is what
  \verb+remove+, \verb+rmflip+ and \verb+exogenise+ act on, and why a \verb+.mod+ file is
  read with its \verb+name+ tags, or paired as described earlier.

\end{frame}


\begin{frame}[c,fragile]
  \frametitle{Tag equations}

  \begin{lstlisting}[style=MatlabConsole]

    >> model.tag('h', 'kind', 'static');
    >> model.tag('k', 'kind', 'euler');
   \end{lstlisting}

   \bigskip

  \begin{itemize}

  \item The first argument specifies the equation's name, the second
    argument denotes the tag's name, and the third argument represents
    the tag's value.\newline

  \item An equation can have any number of tags added to it.\newline

  \item A tag can be named anything, except for 'name', which is
    specifically reserved for linking equations to an endogenous
    variable.\newline

  \end{itemize}

\end{frame}


\begin{frame}[fragile]
  \frametitle{Tag equations -- an example --}

  Two sectors, \verb+m+ and \verb+s+. Tags are set with implicit loops, and a tag
  value may use the index:

  \begin{lstlisting}[style=MatlabConsole,basicstyle=\ttfamily\scriptsize]
m = modBuilder();
m.add('Y_$1', 'Y_$1 = A_$1*K_$1(-1)^alpha*L_$1^(1-alpha)', ...
      {'m', 's'});
m.add('K_$1', 'K_$1 = (1-delta)*K_$1(-1) + I_$1', {'m', 's'});
m.add('C', 'C = Y_m + Y_s - I_m - I_s');
% parameters and exogenous variables declared as usual
m.tag('Y_$1', 'type', 'production', {'m', 's'});
m.tag('K_$1', 'type', 'accumulation', {'m', 's'});
m.tag('Y_$1', 'sector', '$1', {'m', 's'});
m.tag('K_$1', 'sector', '$1', {'m', 's'});
m.tag('C', 'type', 'aggregation');
  \end{lstlisting}

  \begin{itemize}
  \item The tags of an equation are the fields of a structure: \verb+m.tags.Y_m+ holds
    \verb+name+ \verb+'Y_m'+, \verb+type+ \verb+'production'+ and \verb+sector+ \verb+'m'+.
  \item The next frames use them to find equations, to extract a block, and to remove
    one.
  \end{itemize}

\end{frame}


\begin{frame}[fragile]
  \frametitle{Tag equations -- find equations --}

  \begin{lstlisting}[style=MatlabConsole,basicstyle=\ttfamily\scriptsize]
>> m.listeqbytag('type', 'production')
   % Y_m Y_s
>> m.listeqbytag('sector', 'm')
   % Y_m K_m
>> m.listeqbytag('sector', 'm', 'type', 'production')
   % Y_m
>> m.listeqbytag('type', 'acc.*')
   % K_m K_s
>> m.listeqbytag('type', 'production|accumulation')
   % Y_m Y_s K_m K_s
  \end{lstlisting}

  \begin{itemize}\small
  \item Several criteria must all hold: the third call keeps only the production
    equation of sector \verb+m+.
  \item A value is a regular expression matched against the whole tag value:
    \verb+'acc.*'+ matches \verb+accumulation+, and \verb+|+ reads as ``or''.
  \item No match is an error, not an empty list.
  \end{itemize}

\end{frame}


\begin{frame}[fragile]
  \frametitle{Tag equations -- select a block --}

  \begin{lstlisting}[style=MatlabConsole,basicstyle=\ttfamily\scriptsize]
>> p = m.select('sector', 's');
>> p = m{bytag('sector', 's')};     % the same
  \end{lstlisting}

  \begin{itemize}
  \item \verb+select+ returns a model made of the equations that match and of the
    symbols they use. Here \verb+p+ holds the equations of \verb+Y_s+ and \verb+K_s+,
    the parameters \verb+A_s+, \verb+alpha+ and \verb+delta+, and the exogenous
    variables \verb+L_s+ and \verb+I_s+.\newline
  \item A variable whose equation is left out becomes exogenous in the selection:
    \verb+m{bytag('type', 'production')}+ keeps the equations of \verb+Y_m+ and
    \verb+Y_s+, with \verb+K_m+ and \verb+K_s+ exogenous.
  \end{itemize}

\end{frame}


\begin{frame}[fragile]
  \frametitle{Tag equations -- remove a block, write the model --}

  \verb+rm+ and \verb+remove+ take the same selector:

  \begin{lstlisting}[style=MatlabConsole,basicstyle=\ttfamily\scriptsize]
>> m.rm(bytag('sector', 'm'));       % equations left: Y_s K_s C
  \end{lstlisting}

  \small As on the frame on removal, \verb+Y_m+, still used by \verb+C+, becomes
  exogenous, while \verb+K_m+ and \verb+A_m+, used only by the removed equations, go.

  \medskip
  \normalsize Tags are written to the \verb+.mod+ file, where Dynare accepts any tag,
  and read back with it:

  \begin{lstlisting}[language=Dynare,basicstyle=\scriptsize\ttfamily,numbers=none]
[name = 'Y_m', sector = 'm', type = 'production']
Y_m = A_m*K_m(-1)^alpha*L_m^(1-alpha);

[name = 'C', type = 'aggregation']
C = Y_m + Y_s - I_m - I_s;
  \end{lstlisting}

  \small So \verb+modBuilder('sectors.mod')+ gives a model on which the same selections
  work, and tags written by hand in a \verb+.mod+ file serve the same purpose.

\end{frame}


\begin{frame}[c,fragile]
  \frametitle{Type symbols}

  \begin{itemize}

  \item \verb+model.symbols+ is the list of unclassified symbols.\newline

  \item This property must be empty to finalize the model.\newline

  \end{itemize}

  \begin{lstlisting}[style=MatlabConsole]

    >> model.parameter('rho', .9, 'texname', '\rho');
    >> model.exogenous('epsilon', .0);
   \end{lstlisting}

   \medskip

  \begin{itemize}

  \item Calibration occurs upon declaring a parameter but can be modified subsequently.\newline

  \item The value assigned to an exogenous variable determines its long-run level.

  \end{itemize}

\end{frame}


\begin{frame}[c,fragile]
  \frametitle{Converting symbols}

  \begin{itemize}

  \item Methods \verb+parameter+ and \verb+exogenous+ can also be used to convert a symbol.\newline

  \item These methods cannot be applied to an endogenous variable.\newline

  \item To exogenize an endogenous variable one can:
    \begin{itemize}
    \item[--] Remove an equation (methods \verb+rm+ or \verb+remove+)
    \item[--] Flip an exogenous variable and and endogenous variable (\verb+flip+ method).
    \end{itemize}

  \end{itemize}

    \bigskip

    \dnote{The method \texttt{endogenous} does not change the type of a variable, but only serves to provide the steady state level of an endogenous variable.}

    \begin{lstlisting}[style=MatlabConsole]

    >> model.endogenous('y', .0);
  \end{lstlisting}

\end{frame}


\begin{frame}[c,fragile]
  \frametitle{Change an equation}

  \begin{itemize}
  \item \verb+change+ method: replace a whole equation.
  \item \verb+subs+ method: replace an expression by another expression. Both arguments are parsed into AST trees, so the replacement is precedence-safe and lag-aware: substituting \verb+mc+ by \verb+w/mpl+ correctly rewrites \verb+mc(-1)+ as \verb+w(-1)/mpl(-1)+.
  \item \verb+substitute+ method: regex-based rewrite on the equation text. Use when the target is a textual pattern (e.g.\ an indexed family) rather than a parsable expression. Warns when the target is a single symbol that \verb+subs+ would handle better.
  \item \verb+rename+ method: rename a symbol in all equations (preserves lag and \verb+STEADY_STATE+ wrapping).
  \end{itemize}

    \medskip

    \dnote{If new symbols are introduced with the changes, they must be typed after.}

    \medskip

    \begin{lstlisting}[style=MatlabConsole]
      >> model.change('y', 'y=rho*y(-1)+phi*y(-2)+e');
      >> model.subs('y(-2)', 'y(-3)', 'y');
  \end{lstlisting}

\end{frame}


\begin{frame}[c,fragile]
  \frametitle{Search for symbols within a model}

  \begin{itemize}

  \item The \verb+lookfor+ method searches for any symbol, such as a parameter or variable, within a model and displays the equations in which the symbol appears.\newline

  \end{itemize}

    \medskip

    \begin{lstlisting}[style=MatlabConsole]

      >> model.lookfor('k')

      Endogenous variable k appears in 3 equations:

      1/beta = ((exp(b)*c)/(exp(b(+1))*c(+1)))*(exp(b(+1))*alpha*y(+1)/k+(1-delta))

      y = exp(a)*(k(-1)^alpha)*(h^(1-alpha))

      k = exp(b)*(y-c)+(1-delta)*k(-1)
  \end{lstlisting}

\end{frame}


\begin{frame}[c,fragile]
  \frametitle{Use matlab to write an equation -- I --}

  \begin{itemize}

  \item The features of the MATLAB language, such as loops and conditional statements, can be employed to programmatically generate equations.\newline

  \item Suppose we want to add the following equation:
    \[
      y_t = \rho_1 y_{t-1} + \rho_2 y_{t-2} + \dots + \rho_{12} y_{t-12} + \varepsilon_t
    \]

    \medskip

  \item Define a general MATLAB routine generating AR equations:\newline

    \lstinputlisting[
frame=none,
numbers=left,
basicstyle=\small,
style=Matlab-editor
]{../models/ar.m}

  \end{itemize}

\end{frame}


\begin{frame}[c,fragile]
  \frametitle{Use matlab to write an equation -- II --}

  \lstinputlisting[
  firstline=1,
  lastline =8,
  basicstyle=\small,
  frame=none,
  numbers=left,
  style=Matlab-editor
  ]{../models/ar12.m}

  \bigskip

  \begin{lstlisting}[style=MatlabConsole]

      >> m.y.equations{2}

      ans =

      'y = rho1*y(-1) + rho2*y(-2) + rho3*y(-3) + rho4*y(-4) + rho5*y(-5) + rho6*y(-6) + rho7*y(-7) + rho8*y(-8) + rho9*y(-9) + rho10*y(-10) + rho11*y(-11) + rho12*y(-12) + e'
  \end{lstlisting}

\end{frame}


\begin{frame}[c,fragile]
  \frametitle{Use matlab to write equations -- III --}

  \lstinputlisting[
  firstline=1,
  lastline =10,
  basicstyle=\tiny,
  frame=none,
  numbers=left,
  style=Matlab-editor
  ]{../models/ar12panel_0.m}

\end{frame}


\begin{frame}[c,fragile]
  \frametitle{Use matlab to write equations -- IV --}

  Same with an implicit loop:

  \bigskip

  \lstinputlisting[
  firstline=1,
  lastline =5,
  basicstyle=\tiny,
  frame=none,
  numbers=left,
  style=Matlab-editor
  ]{../models/ar12panel_1.m}

\end{frame}


\begin{frame}[c,fragile]
  \frametitle{Use matlab to write equations -- V --}

  Implicit loops can be nested:

  \bigskip

  \lstinputlisting[
  firstline=1,
  lastline =6,
  basicstyle=\tiny,
  frame=none,
  numbers=left,
  style=Matlab-editor
  ]{../models/smc.m}

  $\Rightarrow$ The model now includes 40 additional equations! However, users still need to manually input the symbols introduced in the loops.\newline

  \begin{lstlisting}[style=MatlabConsole]
      >> m.Y_FR_1.equations{2}
      ans =
      'Y_FR_1 = A_FR*K_FR_1^alpha*L_FR_1^(1-alpha)'
  \end{lstlisting}

\end{frame}


\begin{frame}[c,fragile]
  \frametitle{Steady state -- the model --}

  \scalebox{.6}{
  \lstinputlisting[
  firstline=1,
  lastline =23,
  basicstyle=\tiny,
  frame=none,
  numbers=left,
  style=Matlab-editor
  ]{../models/rbc13.m}}

\end{frame}


%==============================================================================
% Steady state: declaring the solution
%==============================================================================

\begin{frame}[c,fragile]
  \frametitle{Steady state -- declare it --}

  When the steady state has a closed form, declare it:

  \begin{itemize}

  \item The \verb+steady+ method attaches a steady-state expression to an
    endogenous variable. The expression may reference parameters and the steady
    state of other variables.\newline

  \item \verb+steady_aux+ declares an intermediate quantity that is not a model
    variable (a ratio, a scale), \verb+steady_call+ registers a call assigning
    several variables at once.\newline

  \item \verb+checksteady+ validates the declarations and returns them in
    dependency order, refusing unknown symbols and circular definitions.\newline

  \item \verb+write+ emits them as a \verb+steady_state_model+ block, so that
    Dynare runs no numerical solve at all.

  \end{itemize}

\end{frame}


\begin{frame}[fragile]
  \frametitle{Steady state -- declare it, the whole block --}

  For the model above, capital, output and consumption per hour follow from the
  Euler equation and the resource constraint, and hours from labour supply:

  \begin{lstlisting}[style=MatlabConsole,basicstyle=\ttfamily\scriptsize]
model.steady('a', '0');
model.steady('b', '0');
model.steady_aux('kh', '(alpha/(1/beta-1+delta))^(1/(1-alpha))');
model.steady_aux('yh', 'kh^alpha');
model.steady_aux('ch', 'yh - delta*kh');
model.steady('h', '((1-alpha)*yh/(theta*ch))^(1/(1+psi))');
model.steady('k', 'kh*h');
model.steady('y', 'yh*h');
model.steady('c', 'ch*h');
model.write('rbc', steady_state_model=true, steady=true);
  \end{lstlisting}

  \begin{itemize}\small
  \item \verb+checksteady+ orders them \verb+a b kh yh ch h k y c+; the ratios stay
    local to the \verb+steady_state_model+ block.
  \item Dynare runs the written file with no numerical solve:
    $h^\star = 0.2918$, $k^\star = 11.08$, $y^\star = 1.081$, $c^\star = 0.8036$.
  \end{itemize}

\end{frame}


%==============================================================================
% Steady state: computing the solution numerically
%==============================================================================

\begin{frame}[c,fragile]
  \frametitle{Steady state -- compute it numerically --}

  When no closed form is at hand, two routes:

  \begin{itemize}

  \item Let Dynare solve the whole static system: give initial guesses, and
    \verb+write+ emits an \verb+initval+ block and the \verb+steady+ command (next
    frame).\newline

  \item Or solve within \verb+modBuilder+, one equation at a time:
    \begin{itemize}
    \item[--] \verb+evaluate+ evaluates a static equation at the current values, and
      returns the left-hand side (\verb+lhs+), the right-hand side (\verb+rhs+) and
      the residual \verb+resid=lhs-rhs+;
    \item[--] \verb+solve+ solves the static version of an equation for one symbol,
      an endogenous variable, an exogenous variable or a parameter.
    \end{itemize}

  \end{itemize}

\end{frame}


\begin{frame}[fragile]
  \frametitle{Steady state -- let Dynare solve it --}

  \begin{lstlisting}[style=MatlabConsole,basicstyle=\ttfamily\scriptsize]
% initial guesses, one per endogenous variable
model.endogenous('a', 0);
model.endogenous('b', 0);
model.endogenous('h', 1/3);
model.endogenous('k', 10);
model.endogenous('y', 1);
model.endogenous('c', 0.8);
model.write('rbc', initval=true, steady=true);
  \end{lstlisting}

  \small writes, after the model block (blank lines dropped):

  \begin{lstlisting}[language=Dynare,basicstyle=\scriptsize\ttfamily,numbers=none]
initval;
  a = 0.000000;
  b = 0.000000;
  y = 1.000000;
  c = 0.800000;
  h = 0.333333;
  k = 10.000000;
end; // initval block
steady;
  \end{lstlisting}

  Dynare's \verb+steady+ solves the whole static system from these guesses, and here
  finds the point declared on the previous frames, to within $4\cdot 10^{-7}$.
  \verb+steady_options+ passes options to its solver, e.g.\ \verb+{'maxit', 100}+.

\end{frame}


\begin{frame}[c,fragile]
  \frametitle{Steady state -- equation by equation --}

  \scalebox{1}{
  \lstinputlisting[
  firstline=27,
  lastline =36,
  basicstyle=\tiny,
  frame=none,
  numbers=left,
  style=Matlab-editor
  ]{../models/rbc13.m}}

  \begin{itemize}\small
  \item Each \verb+solve+ handles one static equation for one symbol; the order of
    the calls does the rest, as by hand.
  \item The last call calibrates: with hours fixed at $1/3$, the equation of
    \verb+h+ is solved for \verb+psi+.
  \end{itemize}

\end{frame}


%==============================================================================
% Steady state: deriving the solution
%==============================================================================

\begin{frame}[c,fragile]
  \frametitle{Steady state plan -- the idea --}

  \begin{itemize}

  \item Writing the steady state by hand is the tedious part of building a
    model. The \verb+steady_plan+ method derives it instead.\newline

  \item Three steps:
    \begin{itemize}
    \item[--] pair each equation with the variable it pins, by bipartite
      matching on the \textbf{static} residuals (option \verb+Match+);
    \item[--] decompose the dependency graph into strongly connected components,
      in topological order;
    \item[--] close each block symbolically: isolation for a single equation,
      elimination or Cramer for a small simultaneous block.
    \end{itemize}

  \item What comes out is a \textbf{cascade of closed forms}, in evaluation
    order, and the list of variables (if any) that resisted.

  \end{itemize}

\end{frame}


\begin{frame}[c,fragile]
  \frametitle{Steady state plan -- a small model --}

  \lstinputlisting[
  style=Matlab-editor,
  firstline=1,
  lastline=18,
  basicstyle=\tiny\mlttfamily,
  frame=none,
  numbers=left
  ]{../models/rbc14.m}

\end{frame}


\begin{frame}[c,fragile]
  \frametitle{Steady state plan -- deriving it --}

  \lstinputlisting[
  style=Matlab-editor,
  firstline=1,
  lastline=9,
  basicstyle=\scriptsize\mlttfamily,
  frame=none,
  numbers=left
  ]{../models/rbc15.m}

  \medskip

  \begin{itemize}

  \item The level of \verb|a| is not declared: the plan closes it from its own
    process, \verb|a = rho*a(-1) + e|, as $a = e/(1-\rho)$.\newline

  \item \verb+PropagateKnown+ inlines what is already known -- here the calibrated
    $e=0$ -- before the recognisers run: $a=0$, and \verb|exp(a)| drops out of the
    closed forms of $k$ and $y$.\newline

  \item \verb+calibrate+ swaps a role: hours are fixed at $1/3$, $\theta$
    becomes an unknown, and gets a closed form where \verb+solve('h', 'psi', 0)+
    searched numerically.

  \end{itemize}

\end{frame}


\begin{frame}[c,fragile]
  \frametitle{Steady state plan -- the plan --}

  \begin{lstlisting}[style=MatlabConsole,basicstyle=\ttfamily\scriptsize]
Steady-state plan: 5 block(s)

  Block 1 [anchor]
    vars: a
    closed form: a = 0

  Block 2 [simultaneous, 2 vars]
    vars: y, k
    depends on (already solved): a
    external constants: h, alpha, delta, beta
    closed form: k = ((-1 + delta + beta ^ -1) / alpha * h ^ (-(1 - alpha))) ^ ((-1 + alpha) ^ -1)
    closed form: y = h ^ (1 - alpha) * k ^ alpha

  Block 3 [trivial]
    vars: i
    closed form: i = -(-k + k * (1 - delta))

  Block 5 [trivial]
    vars: theta
    closed form: theta = y * (1 - alpha) / c * h ^ (-1 - psi)
  \end{lstlisting}

  $\Rightarrow$ The calibration of $\theta$ consistent with $h^{\star}=1/3$ is
  derived, not searched for.

\end{frame}


\begin{frame}[c,fragile]
  \frametitle{Steady state plan -- options --}

  \begin{itemize}

  \item \verb+Match+ recomputes the equation/variable pairing instead of using
    the declaration order. A large simultaneous block usually falls apart into
    singletons.\newline

  \item \verb+Anchors+ lists the variables whose level the modeller fixes -- a
    normalisation, a scale, a calibration constant. Each anchor frees one
    equation, which becomes a consistency condition.\newline

  \item \verb+PropagateKnown+ inlines every value already known, including the
    calibrated levels of the exogenous processes.\newline

  \item \verb+MaxBlockSize+ caps the size of the blocks the symbolic machinery
    attempts. Elimination cost grows quickly with block size.\newline

  \item \verb+suggest_anchors+ and \verb+suggest_calibrations+ propose the
    anchors, or the role swaps, that would close what resisted.

  \end{itemize}

\end{frame}


\begin{frame}[c,fragile]
  \frametitle{Steady state plan -- writing it out --}

  \verb+apply_steady_plan+ turns the plan into steady-state declarations, which
  \verb+write+ emits as an analytical block:

  \bigskip

  \begin{lstlisting}[language=Dynare,basicstyle=\footnotesize\ttfamily,numbers=none]
steady_state_model;

	h = 0.333333333333333;
	a = 0;
	k = ((-1 + delta + beta ^ -1) / alpha * h ^ (-(1 - alpha))) ^ ((-1 + alpha) ^ -1);
	y = h ^ (1 - alpha) * k ^ alpha;
	i = -(-k + k * (1 - delta));
	c = -(i - y);
	theta = y * (1 - alpha) / c * h ^ (-1 - psi);

end;
  \end{lstlisting}

\end{frame}


\begin{frame}[c,fragile]
  \frametitle{Steady state plan -- blocks that resist --}

  \begin{itemize}

  \item Not every steady state closes in closed form. When an elimination
    stalls, the plan reports the \textbf{reduced system} -- the open variables
    and their residuals, with every closed form substituted in.\newline

  \item \verb+apply_steady_plan(GenerateHelpers=true)+ then writes that system
    to a self-contained MATLAB routine, called from the
    \verb+steady_state_model+ block: analytic prologue, small numerical solve,
    analytic epilogue.\newline

  \item The generated routine differentiates by \textbf{complex step},
    $f'(x)=\Im\left(f(x+\imath h)\right)/h$ with $h=10^{-20}$, exact to machine
    precision. Option \verb+Jacobian+ selects \verb+auto+, \verb+complex-step+
    or \verb+symbolic+.\newline

  \item Option \verb+Solver+ selects the damped Newton written into the file
    (\verb+newton+) or the solvers shipped with Dynare (\verb+dynare+).

  \end{itemize}

\end{frame}


\begin{frame}[c,fragile]
  \frametitle{Steady state plan -- homogeneous blocks --}

  \begin{itemize}

  \item A block can pin the \textbf{ratios} of its variables without pinning
    their level: the Calvo price and wage recursions are the textbook case.\newline

  \item Solving it as it stands would land on the trivial all-zero point.
    Instead the block is re-parametrised in one gauge, and the gauge is isolated
    from a consistency equation left over by the anchors.\newline

  \item This is the hand technique of solving in ratios and backing the level
    out of a discarded equation -- aggregate production for output,
    labour-market clearing for the wage recursions.\newline

  \item Reported in the plan as \verb+backout+, with the gauge and the equation
    that pinned it.

  \end{itemize}

\end{frame}


%==============================================================================
% LaTeX output
%==============================================================================

\begin{frame}[c,fragile]
  \frametitle{LaTeX -- the model --}

  \begin{lstlisting}[style=MatlabConsole]

    >> model.tex_model('rbcmodel.tex');
  \end{lstlisting}

  \medskip

  Uses the \verb+texname+ metadata declared with each symbol:

  \medskip

  \scalebox{0.85}{$
    \begin{aligned}
      a_t &= \rho\,a_{t-1} + \varepsilon_t \\
      y_t &= e^{a_t}\,k_{t-1}^{\alpha}\,h_t^{1 - \alpha} \\
      k_t &= \left(1 - \delta\right)\,k_{t-1} + i_t \\
      y_t &= c_t + i_t \\
      \frac{1}{\beta} &= \frac{c_t}{c_{t+1}}\,\left(\frac{\alpha\,y_{t+1}}{k_t} + 1 - \delta\right) \\
      \theta\,c_t\,h_t^{\psi} &= \frac{\left(1 - \alpha\right)\,y_t}{h_t}
    \end{aligned}$}

\end{frame}


\begin{frame}[c,fragile]
  \frametitle{LaTeX -- the steady state --}

  \begin{itemize}

  \item \verb+tex_steady_state_system+ renders the static system, every level
    starred.\newline

  \item \verb+tex_steady_state_plan+ renders the derivation itself: the blocks,
    in order, with what each one closes.\newline

  \item \verb+tex_linearise+ renders the model linearised around the steady
    state.

  \end{itemize}

  \medskip

  \scalebox{0.85}{$
    \begin{aligned}
      &y^{\star} - e^{a^{\star}}\,\left. h^{\star} \right.^{1 - \alpha}\,\left. k^{\star} \right.^{\alpha} = 0 \\
      &\delta - 1 - \frac{\alpha\,y^{\star}}{k^{\star}} + \frac{1}{\beta} = 0 \\
      &c^{\star}\,\theta\,\left. h^{\star} \right.^{\psi} - \frac{y^{\star}\,\left(1 - \alpha\right)}{h^{\star}} = 0
    \end{aligned}$}

\end{frame}


\begin{frame}[c,fragile]
  \frametitle{LaTeX -- the linearised model --}

  \begin{lstlisting}[style=MatlabConsole]

    >> vars = {'a','y','k','i','c','h'};
    >> model.tex_linearise(vars, LevelVars={'a'}, Evaluate=true);
  \end{lstlisting}

  \medskip

  \scalebox{0.75}{$
    \begin{aligned}
      \left(a_t - a^{\star}\right) &= 0.95\,\left(a_{t-1} - a^{\star}\right) \\
      \hat{y}_t &= \left(a_t - a^{\star}\right) + 0.36\,\hat{k}_{t-1} + 0.64\,\hat{h}_t \\
      \hat{k}_t &= 0.975\,\hat{k}_{t-1} + 0.025\,\hat{i}_t \\
      \hat{i}_t &= 3.9001\,\hat{y}_t - 2.9001\,\hat{c}_t \\
      \hat{c}_t &= -0.03475\,\hat{y}_{t+1} + 0.03475\,\hat{k}_t + \hat{c}_{t+1} \\
      \hat{h}_t &= 0.5\,\hat{y}_t - 0.5\,\hat{c}_t
    \end{aligned}$}

  \medskip

  \begin{itemize}

  \item Each row is normalised on the variable its equation pins, so the
    coefficients read as elasticities.\newline

  \item \verb+LevelVars+ marks the variables entering as $x_t-x^{\star}$ rather
    than $\hat{x}_t$, for a zero or negative steady state. \verb+Evaluate=false+
    keeps the coefficients symbolic: $\hat{k}_t =
    \left(1-\delta\right)\hat{k}_{t-1} +
    \frac{i^{\star}}{k^{\star}}\,\hat{i}_t$.

  \end{itemize}

\end{frame}


\begin{frame}[c,fragile]
  \frametitle{LaTeX -- the linearised model, symbolically --}

  \begin{lstlisting}[style=MatlabConsole]

    >> model.tex_linearise(vars, LevelVars={'a'});
  \end{lstlisting}

  \medskip

  \scalebox{0.72}{$
    \begin{aligned}
      \left(a_t - a^{\star}\right) &= \rho\,\left(a_{t-1} - a^{\star}\right) \\
      \hat{y}_t &= \frac{e^{a^{\star}}\,\left. h^{\star} \right.^{1 - \alpha}\,\left. k^{\star} \right.^{\alpha}}{y^{\star}}\,\left(a_t - a^{\star}\right)
                 + \frac{\alpha\,e^{a^{\star}}\,\left. h^{\star} \right.^{1 - \alpha}\,\left. k^{\star} \right.^{\alpha}}{y^{\star}}\,\hat{k}_{t-1}
                 + \frac{e^{a^{\star}}\,\left(1 - \alpha\right)\,\left. h^{\star} \right.^{1 - \alpha}\,\left. k^{\star} \right.^{\alpha}}{y^{\star}}\,\hat{h}_t \\
      \hat{k}_t &= \left(1 - \delta\right)\,\hat{k}_{t-1} + \frac{i^{\star}}{k^{\star}}\,\hat{i}_t \qquad
      \hat{i}_t = \frac{y^{\star}}{i^{\star}}\,\hat{y}_t - \frac{c^{\star}}{i^{\star}}\,\hat{c}_t \\
      \hat{c}_t &= -\frac{\alpha\,y^{\star}}{k^{\star}\,\left(1 - \delta + \frac{\alpha\,y^{\star}}{k^{\star}}\right)}\,\hat{y}_{t+1}
                 + \frac{\alpha\,y^{\star}}{k^{\star}\,\left(1 - \delta + \frac{\alpha\,y^{\star}}{k^{\star}}\right)}\,\hat{k}_t + \hat{c}_{t+1} \\
      \hat{h}_t &= \frac{y^{\star}\,\left(1 - \alpha\right)\,\left. h^{\star} \right.^{-2}}{y^{\star}\,\left(1 - \alpha\right)\,\left. h^{\star} \right.^{-2} + c^{\star}\,\psi\,\theta\,\left. h^{\star} \right.^{\psi - 1}}\,\hat{y}_t
                 - \frac{c^{\star}\,\theta\,\left. h^{\star} \right.^{\psi - 1}}{y^{\star}\,\left(1 - \alpha\right)\,\left. h^{\star} \right.^{-2} + c^{\star}\,\psi\,\theta\,\left. h^{\star} \right.^{\psi - 1}}\,\hat{c}_t
    \end{aligned}$}

  \medskip

  $\Rightarrow$ Exact, but written in steady-state levels: the coefficients of
  the $\hat{y}_t$ row are $1$, $\alpha$ and $1-\alpha$, and nothing says so.

\end{frame}


\begin{frame}[c,fragile]
  \frametitle{LaTeX -- elasticities in the deep parameters --}

  \verb+Substitute=true+ inlines the declared steady state into each coefficient
  and reduces it:

  \begin{lstlisting}[style=MatlabConsole]

    >> model.tex_linearise(vars, LevelVars={'a'}, Substitute=true);
  \end{lstlisting}

  \medskip

  \scalebox{0.78}{$
    \begin{aligned}
      \left(a_t - a^{\star}\right) &= \rho\,\left(a_{t-1} - a^{\star}\right) \\
      \hat{y}_t &= \left(a_t - a^{\star}\right) + \alpha\,\hat{k}_{t-1} + \left(1 - \alpha\right)\,\hat{h}_t \\
      \hat{k}_t &= \left(1 - \delta\right)\,\hat{k}_{t-1} + \delta\,\hat{i}_t \\
      \hat{i}_t &= \frac{1 - \beta + \beta\,\delta}{\alpha\,\beta\,\delta}\,\hat{y}_t
                 - \frac{1 - \beta - \alpha\,\beta\,\delta + \beta\,\delta}{\alpha\,\beta\,\delta}\,\hat{c}_t \\
      \hat{c}_t &= -\left(1 - \beta + \beta\,\delta\right)\,\hat{y}_{t+1} + \left(1 - \beta + \beta\,\delta\right)\,\hat{k}_t + \hat{c}_{t+1} \\
      \hat{h}_t &= \frac{1}{1 + \psi}\,\hat{y}_t - \frac{1}{1 + \psi}\,\hat{c}_t
    \end{aligned}$}

  \medskip

  $\Rightarrow$ The investment share is $\delta$, the Euler coefficient
  $1-\beta\left(1-\delta\right)$, the labour elasticity $1/(1+\psi)$: the
  textbook coefficients, derived. Among the algebraically equal forms of each
  coefficient, the one printed is the one \verb+ast.complexity+ scores as the
  simplest to read.

\end{frame}


%==============================================================================
% Examples
%==============================================================================

\begin{frame}[c,fragile]
  \frametitle{Examples -- Smets-Wouters --}

  \begin{itemize}

  \item \verb+examples/sw+ builds the nonlinear Smets-Wouters model: 78
    equations, 78 endogenous variables, 17 shock processes generated rather than
    typed.\newline

  \item \verb+sw.m+ writes the steady state by hand. Its companion
    \verb+sw_from_anchors.m+ derives it instead, from \textbf{seven} declared
    anchors -- two normalisations, the efficient Tobin's $Q$, capacity
    utilisation and its efficient twin, the inverse markup, the hours
    scale.\newline

  \item The other 54 expressions are derived, gauge backout included; the 17
    shock levels are identified without being declared.\newline

  \item The script then executes the emitted cascade and checks every equation
    at the recovered point.

  \end{itemize}

\end{frame}


\begin{frame}[c,fragile]
  \frametitle{Examples -- two countries --}

  \begin{itemize}

  \item \verb+examples/two-country+ is the other end: two RBC economies with
    perfect risk sharing, whose steady state is only \textbf{conditionally}
    analytical.\newline

  \item The technology processes close, the coupled world core does not: the
    world resource constraint reduces to a sum of powers whose composite
    exponents defeat the recognisers.\newline

  \item One variable is left open and handed to a generated routine; the other
    nine are closed conditional on it. Dynare then computes the steady state
    without a global numerical solve.\newline

  \item Both examples carry a \verb+README+, and neither commits a generated
    artefact: running them regenerates everything.

  \end{itemize}

\end{frame}


\end{document}


% Local Variables:
% ispell-check-comments: exclusive
% ispell-local-dictionary: "american-insane"
% TeX-master: t
% End:
